\section{Method}
\label{sec:method}

We present ApexID, a multi-species animal re-identification model trained to
produce a shared identity-discriminative embedding across heterogeneous
wildlife datasets.

\subsection{Problem Formulation}
\label{sec:method:formulation}

We consider multi-species animal re-identification from image
observations. Let
$\mathcal{D}_{\mathrm{tr}}=\{(x_i,y_i,s_i)\}_{i=1}^{N}$
denote the training set, where $x_i$ is an animal image,
$y_i\in\{1,\ldots,C\}$ is its individual identity label, and
$s_i\in\{1,\ldots,S\}$ is its observed species or coarser animal
category. Each training identity has one parent category,
specified by
$h:\{1,\ldots,C\}\rightarrow\{1,\ldots,S\}$,
with $s_i=h(y_i)$. For category $s$, we define
\begin{equation}
\mathcal{I}_s=\{j:h(j)=s\},\qquad n_s=|\mathcal{I}_s|.
\label{eq:formulation:hierarchy}
\end{equation}
This two-level category--identity hierarchy is constructed
from training annotations and does not require a complete
biological taxonomy.

Our goal is to learn a shared embedding function
$f_\theta:\mathcal{X}\rightarrow\mathbb{R}^{D}$,
where $\mathcal{X}$ is the image input space.
We use its unit-normalized representation
\begin{equation}
\mathbf{z}(x)=\frac{f_\theta(x)}{\|f_\theta(x)\|_2},
\qquad \mathbf{z}_i=\mathbf{z}(x_i).
\label{eq:formulation:embedding}
\end{equation}
Given a query image $x_q$ and a gallery
$\mathcal{G}=\{(x_g,y_g)\}_{g=1}^{M}$, image similarity is
\begin{equation}
\operatorname{sim}(x_q,x_g)
=\mathbf{z}(x_q)^\top\mathbf{z}(x_g).
\label{eq:formulation:similarity}
\end{equation}
For queries represented in the gallery, the retrieval
objective is to rank observations of the same individual
ahead of observations of different individuals, including
those within the same category. Category labels provide
training supervision; retrieval uses image embeddings
without requiring query category labels or category-based
filtering.

We consider both closed-set retrieval and open-world
recognition. Writing
$\mathcal{Y}_{\mathcal{G}}
=\{y_g:(x_g,y_g)\in\mathcal{G}\}$,
closed-set queries belong to $\mathcal{Y}_{\mathcal{G}}$,
whereas open-world queries may also contain identities
absent from the initial gallery. The latter setting
requires recognizing enrolled individuals, detecting novel
identities, and associating repeated observations of
previously unseen individuals. Here, novelty is defined
relative to the initial gallery, independently of whether
the identity appeared in $\mathcal{D}_{\mathrm{tr}}$.
In causal online evaluation, assignments use only the
initial gallery and current and past observations, with
discovered identities retained for subsequent association.

\subsection{Model Architecture}
\label{sec:method:architecture}

ApexID uses ConvNeXt-Small~\cite{liu2022convnet} as its image encoder. We
choose a convolutional backbone as a practical compromise between recognition
capacity and deployment cost. Its local, translation-equivariant processing is
well suited to fine-grained markings and textures that may appear at different
positions in an animal crop, while its hierarchical feature maps capture these
cues across multiple spatial scales. ConvNeXt incorporates design advances
associated with modern vision transformers while retaining efficient
convolutional operations. The Small variant also provides a comparatively
lightweight alternative to large transformer encoders and avoids the
quadratic spatial cost of global self-attention, making training and inference
more practical for a general-purpose Re-ID model.

Given a segmented animal image, the backbone produces a single
$D$-dimensional representation, which is $\ell_2$-normalized to obtain the
retrieval embedding $\mathbf{z}$ defined in
Eq.~\eqref{eq:formulation:embedding}. The encoder is shared across all species;
we do not use species-specific backbone branches or require a species label at
inference.

During training, the shared embedding is connected to two supervision heads.
The identity head contains $K$ learnable sub-centers for every training
identity. The same head is used by both the global identity objective and the
conditional within-category identity objective. The category head contains one
learnable proxy for each observed species or coarser animal category and
provides the category-level term of the hierarchical objective. These heads
affect the organization of the shared embedding space but are not part of the
inference model. At test time, both heads are discarded, and the
ConvNeXt-Small embedding is compared directly using cosine similarity.

\subsection{Hierarchical Loss}
\label{sec:method:loss}

We train ApexID with Balanced Hierarchical Category--Animal identity
Sub-center SphereFace2 (BHCA-SCSF2), abbreviated SCSF2.
The objective combines global identity discrimination with balanced category
supervision and conditional within-category identity discrimination.
The hierarchy contains two supervised levels: the available species or
coarser animal category, and the individual identity; it does not require a
complete biological taxonomy.

Using the notation in Sec.~\ref{sec:method:formulation},
consider a mini-batch of $B$ training observations.

\paragraph{Shared identity sub-centers.}
Following sub-center proxy learning~\cite{deng2020subcenter}, each identity
$j$ has $K$ learnable proxies $\mathbf{w}_{j,k}\in\mathbb{R}^{D}$.
Writing hats for $\ell_2$ normalization, its similarity is
\begin{equation}
c_{ij}=\max_{1\leq k\leq K}\mathbf{z}_i^\top\widehat{\mathbf{w}}_{j,k}.
\label{eq:scsf2:subcenter}
\end{equation}
A positive sample can match any sub-center of its identity, accommodating
multiple appearance modes without assigning explicit mode labels.
For a negative identity, the maximum selects its most similar sub-center.
Global and conditional identity losses share this single proxy bank.

\paragraph{Global identity discrimination.}
We build on the one-versus-all SphereFace2
objective~\cite{wen2022sphereface2}, using its additive angular-margin
variant. Let $\theta_{ij}=\arccos(c_{ij})$ and
\begin{align}
g_t(u)&=2\left(\frac{u+1}{2}\right)^t-1,\\
d_{ij}&=g_t\!\left(\cos\!\left[
\min\!\left(\pi,\theta_{ij}
+m_{\mathrm{ID}}\mathbf{1}_{j=y_i}\right)\right]\right).
\label{eq:scsf2:transform}
\end{align}
Inverse-cosine inputs and angles are clipped for numerical stability.
The implementation uses the detached correction
\begin{equation}
a_{ij}=r\left(c_{ij}+
\operatorname{sg}[d_{ij}-c_{ij}]\right)+b,
\label{eq:scsf2:logits}
\end{equation}
where $\operatorname{sg}$ is stop-gradient, $r>0$ is a scale, and $b$ is
one learnable bias shared across identities and both identity objectives.
Thus the forward logit equals $r d_{ij}+b$, while its implemented
derivative with respect to $c_{ij}$ is $r$.

Define $\ell_i^+=\log(1+\exp(-a_{i y_i}))$ and
$\ell_{ij}^-=\log(1+\exp(a_{ij}))$.
The global objective is
\begin{equation}
\mathcal{L}_{\mathrm{G}}
=\frac{\eta}{rB}\sum_{i=1}^{B}
\left[
\alpha\ell_i^+
+(1-\alpha)\sum_{j\neq y_i}\ell_{ij}^-
\right],
\label{eq:scsf2:global}
\end{equation}
where $\alpha\in(0,1)$ balances the positive term against each negative
term and $\eta>0$ scales the identity objectives.
All other training identities are negatives, including identities in
other categories. The negative terms are summed, not averaged.

\paragraph{Balanced category supervision.}
Each category has one learnable unit-normalized proxy
$\widehat{\mathbf{v}}_s$, with
$q_{is}=\mathbf{z}_i^\top\widehat{\mathbf{v}}_s$.
We apply an ArcFace-style angular margin~\cite{deng2019arcface} to the
target category:
\begin{equation}
u_{is}=
\begin{cases}
\gamma\,\phi_{m_{\mathrm{SP}}}(q_{is}), & s=s_i,\\
\gamma\,q_{is}, & s\neq s_i,
\end{cases}
\label{eq:scsf2:specieslogits}
\end{equation}
where $\phi_m(q)=\cos(\arccos(q)+m)$ with the standard monotonic
continuation near $\pi$, and $\gamma>0$ is the category scale.
Under identity-uniform sampling, categories with more training identities
are encountered more frequently. We moderate this imbalance using
\begin{equation}
\omega_s=
\frac{n_s^{-1/2}}{S^{-1}\sum_{v=1}^{S}n_v^{-1/2}}.
\label{eq:scsf2:weights}
\end{equation}
For anchors $\mathcal{V}$ with valid category labels, the weighted
cross-entropy is
\begin{equation}
\mathcal{L}_{\mathrm{SP}}
=
-\frac{\displaystyle\sum_{i\in\mathcal{V}}\omega_{s_i}
\log\frac{\exp(u_{i s_i})}{\sum_{s=1}^{S}\exp(u_{is})}}
{\displaystyle\sum_{i\in\mathcal{V}}\omega_{s_i}}.
\label{eq:scsf2:species}
\end{equation}
The weights use identity counts rather than image counts and provide
partial compensation for category imbalance.

\paragraph{Conditional identity discrimination.}
To add supervision for distinguishing individuals within a category,
we restrict the negative set to $\mathcal{I}_{s_i}\setminus\{y_i\}$.
This set includes every training identity in that category, irrespective
of its presence in the mini-batch.
For $\mathcal{V}_2=\{i\in\mathcal{V}:n_{s_i}>1\}$, we use
\begin{equation}
\begin{split}
\mathcal{L}_{\mathrm{W}}
=\frac{\eta}{r|\mathcal{V}_2|}\sum_{i\in\mathcal{V}_2}
\Bigg[
&\alpha\ell_i^+\\[-2pt]
&+\frac{1-\alpha}{n_{s_i}-1}
\sum_{\substack{j\in\mathcal{I}_{s_i}\\j\neq y_i}}
\ell_{ij}^-
\Bigg].
\end{split}
\label{eq:scsf2:within}
\end{equation}
The logits, sub-centers, margin, and bias are shared with
Eq.~\eqref{eq:scsf2:global}; no separate conditional identity classifier
is introduced. Averaging sibling negatives makes the additional negative
term depend on their mean error rather than their number.

\paragraph{Combined objective.}
The complete loss is
\begin{equation}
\boxed{
\mathcal{L}_{\mathrm{SCSF2}}
=\mathcal{L}_{\mathrm{G}}
+\frac{\lambda_{\mathrm{H}}}{2}
\left(\mathcal{L}_{\mathrm{SP}}+\mathcal{L}_{\mathrm{W}}\right).
}
\label{eq:scsf2:total}
\end{equation}
Global discrimination retains cross-category negatives, category
supervision encourages coarse organization, and the conditional term adds
supervision among sibling identities.
An empty valid set contributes zero; missing category labels retain the
global identity term, and singleton categories skip only the conditional
term. The released configuration uses $K=2$, $m_{\mathrm{ID}}=0.30$,
$m_{\mathrm{SP}}=0.10$ radians, $r=\gamma=40$, $t=3$, $\alpha=0.9$,
$\eta=20$, and $\lambda_{\mathrm{H}}=0.10$.
At inference, the training heads are discarded and retrieval uses cosine
similarity between normalized image embeddings without category gating.


\subsection{Training Objective and Optimization}
\label{sec:method:training}

\TODO{Give the complete training objective and report the sampling strategy,
optimizer, learning-rate schedule, augmentation, image resolution, batch
composition, and other details needed to reproduce training.}